%LTeX: language=fr-FR
\documentclass[a4paper,11pt,fleqn]{article}

\usepackage{tasks}
%%% COULEURS %%%
\usepackage[table,svgnames]{xcolor}
\definecolor{nombres}{cmyk}{0,.8,.95,0}
\definecolor{gestion}{cmyk}{.75,1,.11,.12}
\definecolor{gestionbis}{cmyk}{.75,1,.11,.12}
\definecolor{grandeurs}{cmyk}{.02,.44,1,0}
\definecolor{geo}{cmyk}{.62,.1,0,0}
\definecolor{algo}{cmyk}{.69,.02,.36,0}
\definecolor{correction}{cmyk}{.63,.23,.93,.06}
\arrayrulecolor{couleur_theme} % Couleur des filets des tableaux

\usepackage[most]{tcolorbox} % Supprimé à cause de conflits avec ProfCollege
% Découpe des cadres d'exercice. Le style est redéfini par la sortie LaTeX
% juste avant un exercice à garder d'un seul tenant, puis rétabli :
%   \tcbset{mathaleaexo/.style={unbreakable}}
% Passer par un style plutôt que par la clé « breakable » directement évite
% de changer la découpe des autres cadres du document (ProfCollege...).
\tcbset{mathaleaexo/.style={breakable}}
\tcbuselibrary{documentation}

%%%%%%%% Modification paramétrage pour footnotes
% Le paquet semble déjà appelé en amont, curieux qu'il n'y ait pas eu de clash
% avec l'appel ci-dessous mais seulement quand j'ai essayé de l'inclure avec l'option lualatex

% \usepackage{hyperref}
% Parametrages
\hypersetup{
    colorlinks=true,% On active la couleur pour les liens. Couleur par défaut rouge
    linkcolor=black,% On définit la couleur pour les liens internes
    % filecolor=magenta,% On définit la couleur pour les liens vers les fichiers locaux      
    urlcolor=black,% On définit la couleur pour les liens vers des sites web
    % pdftitle={Puissance Quatre},% On définit un titre pour le document pdf
    % pdfpagemode=FullScreen,% On fixe l'affichage par défaut à plein écran
    }
% Pour avoir les numérotations arabic y compris dans les environnements tcolorbox
\renewcommand{\thempfootnote}{\arabic{mpfootnote}}
%%%%%%%%

\usepackage{multicol} 
\usepackage{calc}
\usepackage{enumitem}
\usepackage{graphicx}
\usepackage{tabularx}
\usepackage{tablvar}

\setlength{\parindent}{0mm}
\renewcommand{\arraystretch}{1.5}
\renewcommand{\labelenumi}{\textbf{\theenumi{}.}}
\renewcommand{\labelenumii}{\textbf{\theenumii{}.}}
\setlength{\fboxsep}{3mm}

\setlength{\headheight}{14.5pt}

\setlength{\premulticols}{0pt}

\newcounter{ExoMA}

\newlength{\parindentMA}%
\setlength{\parindentMA}{\parindent}
\addtolength{\parindentMA}{30pt}

\usepackage{simplekv}

\setKVdefault[Theme]{Coopmath,Classiques=false,CANs=false}
\defKV[Theme]{Classique=\setKV[Theme]{Coopmath=false}\setKV[Theme]{Classiques}}
\defKV[Theme]{CAN=\setKV[Theme]{Coopmath=false}\setKV[Theme]{CANs}}

\usepackage{fancyhdr}
\fancyhead[C]{}
\fancyfoot{}

\def\bla{}

\NewDocumentCommand\Theme{o m m m m}{%
  \useKVdefault[Theme]%
  \setKV[Theme]{#1}%
  \ifboolKV[Theme]{CANs}{%
    \usepackage[left=1.5cm,right=1.5cm,top=2cm,bottom=2cm]{geometry}%
    \fancypagestyle{premierePage}{%
      \fancyhead[C]{\textsc{\ifx\bla#2\bla Course aux nombres\else#2\fi}}%
      \fancyfoot[R]{\scriptsize Coopmaths.fr -- CC-BY-SA}%
      \setlength{\headheight}{14.5pt}%
      \NewDocumentCommand\dureeCan{}{\ifx\bla#4\bla 9 minutes\else #4\fi}
      \NewDocumentCommand\titreSujetCan{}{\ifx\bla#3\bla\else\textbf{#3}\fi}
    }%
    \pagestyle{premierePage}
    \colorlet{couleur_theme}{black}%
    \colorlet{couleur_numerotation}{couleur_theme}%
    \let\EXO\EXOCan\let\endEXO\endEXOCan%
  }{%
    \renewcommand{\headrulewidth}{0pt}
    \renewcommand{\footrulewidth}{0pt}
    \pagestyle{fancy}
    \ifboolKV[Theme]{Classiques}{%
      \usepackage[left=0.65cm,right=0.65cm,top=2cm,bottom=1.5cm]{geometry}
      \let\EXO\EXOlibre\let\endEXO\endEXOlibre%
      \colorlet{couleur_theme}{black}%
      \colorlet{couleur_numerotation}{couleur_theme}%
      \fancyhead[C]{\bfseries #3}
      \fancyhead[L]{\bfseries #4}
      \fancyfoot[R]{#5}
    }{%
      \usepackage[left=1.5cm,right=1.5cm,top=4cm,bottom=2cm]{geometry}
      \theme{#2}{#3}{#4}{#5}%
      \let\EXO\EXOcoop\let\endEXO\endEXOcoop%
    }%
  }%
}%

\NewDocumentEnvironment{EXOCan}{m m}{%
}{}

\NewDocumentEnvironment{EXOcoop}{m m}{%
  \begin{tcolorbox}[%
    enhanced,
    mathaleaexo,
    colback=white,
    colframe=white,
    coltitle=black,
    title=\IfNoValueTF{#1}{}{\textmd{#1}},%
    overlay unbroken and first={%
      \IfNoValueTF{#2}{}{%
        \node[
        fill=white,
        anchor=east,
        xshift=10pt,
        text=black
        ]
        at (frame.north east)
        {\scriptsize #2};
      }
      \node[anchor=west,xshift=-20pt,yshift=-15pt] 
      at (frame.north west) {\stepcounter{ExoMA}\numb{\arabic{ExoMA}}};
    }
  ]
}{%
  \end{tcolorbox}
}

\NewDocumentEnvironment{EXOlibre}{m m}{%
  \begin{tcolorbox}[%
    enhanced,
    mathaleaexo,
    colback=white,
    colframe=white,
    coltitle=black,
    title=\stepcounter{ExoMA}\textbf{Exercice \arabic{ExoMA}},%
  ]
  \IfNoValueTF{#1}{}{#1\par}%
}{%
  \end{tcolorbox}%
}


%%% MISE EN PAGE %%%
\usepackage{setspace}
\usepackage{pgf}%
\usepackage{pgfplots}
\pgfplotsset{compat=1.18}
\usetikzlibrary{babel,arrows, arrows.meta,calc,fit,patterns,plotmarks,shapes.geometric,shapes.misc,shapes.symbols,shapes.arrows,shapes.callouts, shapes.multipart, shapes.gates.logic.US,shapes.gates.logic.IEC, er, automata,backgrounds,chains,topaths,trees,petri,mindmap,matrix, calendar, folding,fadings,through,positioning,scopes,decorations.fractals,decorations.shapes,decorations.text,decorations.pathmorphing,decorations.pathreplacing,decorations.footprints,decorations.markings,shadows}

\newcommand\numb[1]{ % Dessin autour du numéro d'exercice
  \begin{tikzpicture}[overlay,scale=.8]%,yshift=-.3cm
    \draw[fill=couleur_numerotation,couleur_numerotation](-.4,-0.1)rectangle($(-0.4,.8)+(2em,0)$);
    \draw (-.4,.8) node[white,anchor=north west,inner sep=0.5pt]{\bfseries EX}; 
    \draw[line width=.05cm,couleur_numerotation,fill=white] (0,0)--(.5,.5)--(1,0)--(.5,-.5)--cycle;
    \draw (0.5,0) node[anchor=center,couleur_numerotation]{\large\bfseries#1};
  \end{tikzpicture}
}

% echelle pour le dé
\def\globalscale {0.04}
% abscisse initiale pour les chevrons
\def\xini {3}

\newcommand\theme[4]{%
  \fancyhead[C]{%
    % Tracé du dé
    \begin{tikzpicture}[y=0.80pt, x=0.80pt, yscale=-\globalscale, xscale=\globalscale,remember picture, overlay,transform canvas={xshift=-0.5\linewidth,yshift=2.7cm},fill=couleur_theme]%,xshift=17cm,yshift=9.5cm
      %%%% Arc supérieur gauche%%%%
      \path[fill](523,1424)..controls(474,1413)and(404,1372)..(362,1333)..controls(322,1295)and(313,1272)..(331,1254)..controls(348,1236)and(369,1245)..(410,1283)..controls(458,1328)and(517,1356)..(575,1362)..controls(635,1368)and(646,1375)..(643,1404)..controls(641,1428)and(641,1428)..(596,1430)..controls(571,1431)and(538,1428)..(523,1424)--cycle;
      %%%% Dé face supérieur%%%%
      \path[fill](512,1272)..controls(490,1260)and(195,878)..(195,861)..controls(195,854)and(198,846)..(202,843)..controls(210,838)and(677,772)..(707,772)..controls(720,772)and(737,781)..(753,796)..controls(792,833)and(1057,1179)..(1057,1193)..controls(1057,1200)and(1053,1209)..(1048,1212)..controls(1038,1220)and(590,1283)..(551,1282)..controls(539,1282)and(521,1278)..(512,1272)--cycle;
      %%%% Dé faces gauche et droite%%%%
      \path[fill](1061,1167)..controls(1050,1158)and(978,1068)..(900,967)..controls(792,829)and(756,777)..(753,756)--(748,729)--(724,745)..controls(704,759)and(660,767)..(456,794)..controls(322,813)and(207,825)..(200,822)..controls(193,820)and(187,812)..(187,804)..controls(188,797)and(229,688)..(279,563)..controls(349,390)and(376,331)..(391,320)..controls(406,309)and(462,299)..(649,273)..controls(780,254)and(897,240)..(907,241)..controls(918,243)and(927,249)..(928,256)..controls(930,264)and(912,315)..(889,372)..controls(866,429)and(848,476)..(849,477)..controls(851,479)and(872,432)..(897,373)..controls(936,276)and(942,266)..(960,266)..controls(975,266)and(999,292)..(1089,408)..controls(1281,654)and(1290,666)..(1290,691)..controls(1290,720)and(1104,1175)..(1090,1180)..controls(1085,1182)and (1071,1176)..(1061,1167)--cycle;
      %%%% Arc inférieur bas%%%%
      \path[fill](1329,861)..controls(1316,848)and(1317,844)..(1339,788)..controls(1364,726)and(1367,654)..(1347,591)..controls(1330,539)and(1338,522)..(1375,526)..controls(1395,528)and(1400,533)..(1412,566)..controls(1432,624)and(1426,760)..(1401,821)..controls(1386,861)and(1380,866)..(1361,868)..controls(1348,870)and(1334,866)..(1329,861)--cycle;
      %%%% Arc inférieur gauche%%%%
      \path[fill](196,373)..controls(181,358)and(186,335)..(213,294)..controls(252,237)and(304,190)..(363,161)..controls(435,124)and(472,127)..(472,170)..controls(472,183)and(462,192)..(414,213)..controls(350,243)and(303,283)..(264,343)..controls(239,383)and(216,393)..(196,373)--cycle;
    \end{tikzpicture}
    \begin{tikzpicture}[line cap=round,line join=round,remember picture, overlay, shift={(current page.north west)},yshift=-8.5cm]
      \fill[fill=couleur_theme] (0,5) rectangle (21,6);
      \fill[fill=couleur_theme] (\xini,6)--(\xini+1.5,6)--(\xini+2.5,7)--(\xini+1.5,8)--(\xini,8)--(\xini+1,7)-- cycle;
      \fill[fill=couleur_theme] (\xini+2,6)--(\xini+2.5,6)--(\xini+3.5,7)--(\xini+2.5,8)--(\xini+2,8)--(\xini+3,7)-- cycle;  
      \fill[fill=couleur_theme] (\xini+3,6)--(\xini+3.5,6)--(\xini+4.5,7)--(\xini+3.5,8)--(\xini+3,8)--(\xini+4,7)-- cycle;   
      \node[color=white] at (10.5,5.5) {\LARGE \bfseries{ \MakeUppercase{#4}}};
    \end{tikzpicture}
    \begin{tikzpicture}[remember picture,overlay]
      \node[anchor=north east,inner sep=0pt] at ($(current page.north east)+(0,-.8cm)$) {};
      \node[anchor=west, fill=white, inner sep = 0pt] at ($(current page.north west)+(0.2,-2.3cm)$) {\footnotesize \bfseries{MathALÉA}};
    \end{tikzpicture}
    \begin{tikzpicture}[remember picture,overlay]
      \node[anchor=north east,inner sep=0pt] at ($(current page.north east)+(0,-.8cm)$) {};
      \node[anchor=east, fill=white] at ($(current page.north east)+(-2,-1.5cm)$) {\Huge \textcolor{couleur_theme}{\bfseries{\#}} \bfseries{#2} \textcolor{couleur_theme}{\bfseries \MakeUppercase{#3}}};
    \end{tikzpicture}
  }%
  \fancyfoot[R]{%
    \begin{tikzpicture}[remember picture,overlay]
      \node[anchor=south east] at ($(current page.south east)+(-2,0.25cm)$) {\scriptsize {\bfseries \href{https://coopmaths.fr/}{Coopmaths.fr} -- \href{http://creativecommons.fr/licences/}{CC-BY-SA}}};
    \end{tikzpicture}
    \begin{tikzpicture}[line cap=round,line join=round,remember picture, overlay,xscale=0.5,yscale=0.5, shift={(current page.south west)},xshift=35.7cm,yshift=-6cm]
      \fill[fill=couleur_theme] (\xini,6)--(\xini+1.5,6)--(\xini+2.5,7)--(\xini+1.5,8)--(\xini,8)--(\xini+1,7)-- cycle;
      \fill[fill=couleur_theme] (\xini+2,6)--(\xini+2.5,6)--(\xini+3.5,7)--(\xini+2.5,8)--(\xini+2,8)--(\xini+3,7)-- cycle;  
      \fill[fill=couleur_theme] (\xini+3,6)--(\xini+3.5,6)--(\xini+4.5,7)--(\xini+3.5,8)--(\xini+3,8)--(\xini+4,7)-- cycle;  
    \end{tikzpicture}
  }
  \fancyfoot[C]{}
  \colorlet{couleur_theme}{#1}
  \colorlet{couleur_numerotation}{couleur_theme}
  \def\iconeobjectif{icone-objectif-#1}
  \def\urliconeomethode{icone-methode-#1}
}%

\newcommand*\tikzfootMA{%
  \ifboolKV[Theme]{CANs}{
    \begin{tikzpicture}[remember picture,overlay,shift=(current page.north west)]
      \begin{scope}[x=\textwidth-\oddsidemargin,y=\textheight+5pt]
        \draw[correction,line width=4pt,dashed,dash pattern= on 10pt off 10pt,shift={(-\oddsidemargin,\topmargin)}] (0,0) rectangle (1,-1);
      \end{scope}
    \end{tikzpicture} 
  }{%
  \ifboolKV[Theme]{Classiques}{%
  \begin{tikzpicture}[remember picture,overlay,shift=(current page.south west)]
    \begin{scope}[x={(current page.south east)},y={(current page.north west)}]
      \draw[correction,line width=4pt,dashed,dash pattern= on 10pt off 10pt] ($(0,0)+(5mm,1cm)$) rectangle ($(1,1)+(-5mm,-1.75cm)$);
    \end{scope}
  \end{tikzpicture} 
  }{%
  \begin{tikzpicture}[remember picture,overlay,shift=(current page.south west)]
    \begin{scope}[x={(current page.south east)},y={(current page.north west)}]
      \draw[correction,line width=4pt,dashed,dash pattern= on 10pt off 10pt] ($(0,0)+(5mm,1.5cm)$) rectangle ($(1,1)+(-5mm,-3.75cm)$);
    \end{scope}
  \end{tikzpicture} 
  }%
  }
}

\newenvironment{Correction}{%
  \setcounter{ExoMA}{0}%
  \cfoot{\tikzfootMA}
}{}%

\newcommand\version[1]{
  \fancyhead[R]{
    \begin{tikzpicture}[remember picture,overlay]
    \node[anchor=north east,inner sep=0pt] at ($(current page.north east)+(-.5,-.5cm)$) {\large \textcolor{couleur_theme}{\bfseries V#1}};
    \end{tikzpicture}
  }
}

\usepackage[french]{babel}
\usepackage{icomma}
% loadPackagesFromContent

\newcounter{customfootnote}
\newcounter{tempreset}%
\newcounter{temp}

\makeatletter
\newcommand{\anote}[1]{%
  \refstepcounter{customfootnote}%
  \textsuperscript{\thecustomfootnote}%
  \expandafter\gdef\csname custom@footnote@\thecustomfootnote\endcsname{#1}%
}

\newcommand{\printcustomnotes}{%
  \par\noindent\rule{\linewidth}{0.4pt}\par
  \setcounter{temp}{1}%
  \whileboolexpr{
    test {\ifcsdef{custom@footnote@\thetemp}}
  }{%
    \noindent\textsuperscript{\thetemp}~\csname custom@footnote@\thetemp\endcsname\par
    \stepcounter{temp}%
  }%
}

\newcommand{\resetcustomnotes}{%
  \setcounter{customfootnote}{0}%
  \setcounter{tempreset}{1}%
  \whileboolexpr{
    test {\ifcsdef{custom@footnote@\thetempreset}}
  }{%
    \global\expandafter\let\csname custom@footnote@\thetempreset\endcsname\@undefined
    \stepcounter{tempreset}%
  }%
}
\makeatother
\usepackage[gen]{eurosym}
\usepackage{cancel}
\usepackage{tikz}
\usetikzlibrary{decorations.pathmorphing}
\usetikzlibrary {decorations.pathreplacing}
\newcommand{\e}{\mathrm{e}}
\usepackage{amsmath}
\usepackage{amssymb}
\usepackage{etoolbox}
\newbool{dys}
\setbool{dys}{false}
% ===== POLICES =====
% Le choix est le même en police standard et en police adaptée aux dys ;
% seuls la taille et l'interligne changent entre les deux.
\newcommand{\choiceFonts}[1]{
\ifstrequal{#1}{helvet}{%
  % Helvetica (police standard historique des fiches)
  \usepackage[scaled=1]{helvet}
}{}
\ifstrequal{#1}{Fira}{%
  % Fira Sans + Fira Math
  % Description : Fira Sans est une police moderne, et Fira Math est une version mathématique compatible qui maintient le style sans-serif.
  % Utilisation : Fonctionne bien avec XeLaTeX et LuaLaTeX.
  \usepackage{fontspec}
  \setmainfont{Fira Sans}
  \setsansfont{Fira Sans}
  \usepackage{unicode-math}
  \setmathfont{Fira Math}
}{}
\ifstrequal{#1}{lmodern}{%
  % Latin Modern Sans
  % Description : Une version sans-serif de la célèbre police Latin Modern, qui est compatible avec mathastext.
  % Utilisation : Fonctionne avec pdfLaTeX, XeLaTeX et LuaLaTeX.
  \usepackage{lmodern}
  \renewcommand{\familydefault}{\sfdefault}
  \usepackage{mathastext}
}{}
\ifstrequal{#1}{tgheros}{%
  % TeX Gyre Heros + mathastext
  % Description : TeX Gyre Heros est une alternative à Helvetica, et le package mathastext permet d'utiliser la même police pour les mathématiques.
  % Utilisation : Fonctionne avec pdfLaTeX, XeLaTeX, ou LuaLaTeX.
  \usepackage{tgheros}
  \renewcommand{\familydefault}{\sfdefault}
  \usepackage{mathastext}
}{}
}
\ifbool{dys}{
% POLICE DYS
\choiceFonts{Fira}

%\usepackage{unicode-math}
%\usepackage{fontspec}
%\setmainfont{TeX Gyre Schola}
%\setsansfont{TeX Gyre Schola}
%\setmathfont{TeX Gyre Schola Math}
%\setmainfont{OpenDyslexic}[Scale=1.0]
%\setsansfont{OpenDyslexic}[Scale=1.0]
%\setmathfont{OpenDyslexic}[Scale=1.0,range=up/{Latin,latin,num}]
\usepackage[fontsize=14]{scrextend}
\usepackage{setspace}
\setstretch{1.7}
% Une valeur d'environ 1.2 à 1.7 est souvent recommandée. Cela permet d'augmenter l'espace entre les lignes tout en offrant un léger espacement entre les mots.
\setlength{\spaceskip}{1.2em}
% Une valeur d'environ 1.2em à 1.5em est couramment conseillée. Cela crée un espace plus ample entre les mots, ce qui peut aider à réduire la fatigue visuelle et à améliorer la fluidité de la lecture.
}{
% POLICE STANDARD
%%% EE (24/04/2026) : Cette modif ci-dessous est nécessaire pour accepter ’ comme apostrophe.
% \usepackage[T1]{fontenc}     % Réservé à pdfLaTeX, à remplacer par fontspec en LuaLaTeX
\choiceFonts{helvet}
\usepackage[fontsize=12]{scrextend}
}


% Résultat final mis en valeur dans la correction
\newcommand{\res}[1]{\ensuremath{{\color[HTML]{F15929}\boldsymbol{#1}}}}


% Encadré de propriété
\newtcolorbox{prop}[1]{colback=white, colframe=couleur_theme, title=\textbf{#1}, fonttitle=\bfseries, left=6pt, right=6pt, top=4pt, bottom=3pt, before skip=6pt, after skip=3pt}
% Rubrique dans un encadré : Définition, Propriété, Méthode, Exemples, Démonstration
\newcommand{\rub}[1]{\par\vspace{0.3em}\noindent{\color{couleur_theme}\bfseries #1.}\enspace}
% Exemple ou démonstration à faire en classe : espace blanc en version élève, solution en version prof
\newcommand{\afaire}[2][2.4em]{\par\smallskip\noindent\parbox[t][#1][t]{\linewidth}{\ifdefined\ELEVE\else\color{blue!70!black}#2\fi}\par}

\usepackage{lastpage}
% Liste à puces dans les encadrés
\newenvironment{liste}{\begin{itemize}[label=\textbullet, itemsep=1pt, topsep=1pt, leftmargin=1.6em]}{\end{itemize}}

\Theme[Classique]{nombres}{}{}{}
\begin{document}
\pagestyle{fancy}
\setlength{\headwidth}{\textwidth}
\fancyhf{}
\fancyhead[L]{\textbf{Fiche calcul 6}}
\fancyhead[C]{Racines carrées}
\fancyhead[R]{Seconde}
\fancyfoot[L]{A. Fernandez}
\fancyfoot[R]{\thepage\ / \pageref*{LastPage}}
\renewcommand{\headrulewidth}{0.4pt}

\begin{center}
	{\Large\bfseries Fiche calcul 6 -- Racines carrées}
\end{center}
\vspace{-0.6cm}

\tcbset{mathaleaexo/.style={breakable}}

\begin{prop}{1. Définition}
\rub{Définition}Pour $a \geqslant 0$, $\sqrt{a}$ est le nombre \textbf{positif} dont le carré est $a$.

\rub{Propriétés}
\begin{liste}
	\item Pour $a \geqslant 0$ : \quad $\left(\sqrt{a}\right)^2 = a$.
	\item Pour $a \geqslant 0$ : \quad $\sqrt{a^2} = a$.
\end{liste}

\rub{À connaître}Carrés parfaits : $1$, $4$, $9$, $16$, $25$, $36$, $49$, $64$, $81$, $100$, $121$, $144$.

\rub{Exemples}
\begin{liste}
	\item Calculer $\sqrt{49}$. \afaire[1.2em]{$\sqrt{49} = 7$ car $7^2 = 49$ et $7 \geqslant 0$.}
	\item Calculer $\sqrt{(-5)^2}$. \afaire[1.2em]{$\sqrt{(-5)^2} = \sqrt{25} = 5$ (et non $-5$ : une racine carrée est positive).}
\end{liste}
\end{prop}

\begin{EXO}{}{}
Calculer sans calculatrice.

\begin{tasks}[label=\textbf{\alph*.}, item-indent=2em, after-item-skip=0.6em](4)
	\task $\sqrt{0{,}09}$
	\task $\sqrt{10^6}$
	\task $\left(-\sqrt{3}\right)^2$
	\task $\sqrt{(-8)^2}$
\end{tasks}
\end{EXO}

\begin{prop}{2. Produit}
\rub{Propriété}Pour $a \geqslant 0$ et $b \geqslant 0$ : \quad $\sqrt{a \times b} = \sqrt{a} \times \sqrt{b}$.

\rub{Démonstration}
\afaire[4.6em]{$\sqrt{a}$ et $\sqrt{b}$ sont positifs, donc $\sqrt{a} \times \sqrt{b}$ est positif.

De plus, $\left(\sqrt{a} \times \sqrt{b}\right)^2 = \left(\sqrt{a}\right)^2 \times \left(\sqrt{b}\right)^2 = a \times b$.

Donc $\sqrt{a} \times \sqrt{b}$ est le nombre positif dont le carré est $ab$ : c'est $\sqrt{ab}$.}

\rub{Définition}\textbf{Simplifier} $\sqrt{n}$, c'est l'écrire $a\sqrt{b}$, avec $a$ et $b$ entiers et $b$ le plus petit possible.

\rub{Méthodes}
\begin{liste}
	\item \textbf{Simplifier} : on fait apparaître le plus grand carré parfait qui divise $n$.
	\item \textbf{Réduire} une somme : on simplifie chaque racine, puis on regroupe (comme $2x + 3x = 5x$).
\end{liste}

\rub{Exemples}
\begin{liste}
	\item Calculer $\sqrt{3} \times \sqrt{12}$. \afaire[1.2em]{$\sqrt{3} \times \sqrt{12} = \sqrt{36} = 6$.}
	\item Simplifier $\sqrt{12}$. \afaire[1.2em]{$\sqrt{12} = \sqrt{4 \times 3} = \sqrt{4} \times \sqrt{3} = 2\sqrt{3}$.}
	\item Réduire $\sqrt{12} + \sqrt{27}$. \afaire[1.2em]{$\sqrt{12} + \sqrt{27} = 2\sqrt{3} + 3\sqrt{3} = 5\sqrt{3}$.}
\end{liste}
\end{prop}

\begin{EXO}{}{}
Calculer.

\begin{tasks}[label=\textbf{\alph*.}, item-indent=2em, after-item-skip=0.6em](4)
	\task $\sqrt{2} \times \sqrt{18}$
	\task $\sqrt{5} \times \sqrt{45}$
	\task $\left(3\sqrt{2}\right)^2$
	\task $2\sqrt{3} \times 5\sqrt{6}$
\end{tasks}
\end{EXO}

\begin{EXO}{}{}
Simplifier.

\begin{tasks}[label=\textbf{\alph*.}, item-indent=2em, after-item-skip=0.6em](5)
	\task $\sqrt{48}$
	\task $\sqrt{75}$
	\task $\sqrt{98}$
	\task $\sqrt{500}$
	\task $\sqrt{288}$
\end{tasks}
\end{EXO}

\tcbset{mathaleaexo/.style={unbreakable}}
\begin{EXO}{}{}
Réduire.

\begin{tasks}[label=\textbf{\alph*.}, item-indent=2em, after-item-skip=0.6em](2)
	\task $A = 2\sqrt{50} - 3\sqrt{8} + \sqrt{18}$
	\task $B = \sqrt{75} - 2\sqrt{12} + \sqrt{3}$
	\task $C = 3\sqrt{20} - \sqrt{45} + \sqrt{80}$
	\task $D = \sqrt{2} \times \sqrt{6} + \sqrt{27}$
\end{tasks}
\end{EXO}
\tcbset{mathaleaexo/.style={breakable}}

\begin{prop}{3. Quotient}
\rub{Propriété}Pour $a \geqslant 0$ et $b > 0$ : \quad $\sqrt{\dfrac{a}{b}} = \dfrac{\sqrt{a}}{\sqrt{b}}$.

\rub{Méthode}\textbf{Rendre le dénominateur entier} : multiplier en haut et en bas par cette racine.

\rub{Exemples}Écrire le plus simplement possible :
\begin{liste}
	\item $\sqrt{\dfrac{9}{16}}$ \afaire[2.6em]{$\sqrt{\dfrac{9}{16}} = \dfrac{\sqrt{9}}{\sqrt{16}} = \dfrac{3}{4}$}
	\item $\dfrac{\sqrt{50}}{\sqrt{2}}$ \afaire[2.6em]{$\dfrac{\sqrt{50}}{\sqrt{2}} = \sqrt{\dfrac{50}{2}} = \sqrt{25} = 5$}
	\item $\dfrac{6}{\sqrt{3}}$ \afaire[2.6em]{$\dfrac{6}{\sqrt{3}} = \dfrac{6 \times \sqrt{3}}{\sqrt{3} \times \sqrt{3}} = \dfrac{6\sqrt{3}}{3} = 2\sqrt{3}$}
\end{liste}
\end{prop}

\begin{EXO}{}{}
Écrire le plus simplement possible (pour e à g, simplifier d'abord chaque racine).

\begin{tasks}[label=\textbf{\alph*.}, item-indent=2em, after-item-skip=0.6em](4)
	\task $\sqrt{\dfrac{25}{49}}$
	\task $\dfrac{\sqrt{72}}{\sqrt{8}}$
	\task $\sqrt{\dfrac{12}{27}}$
	\task $\sqrt{\dfrac{8}{50}}$
\end{tasks}
\begin{tasks}[label=\textbf{\alph*.}, item-indent=2em, after-item-skip=0.6em, resume](3)
	\task $\dfrac{\sqrt{50} + \sqrt{8}}{\sqrt{2}}$
	\task $\dfrac{\sqrt{12} - \sqrt{3}}{\sqrt{27}}$
	\task $\dfrac{2\sqrt{20}}{\sqrt{45}}$
\end{tasks}
\end{EXO}

\begin{EXO}{}{}
Écrire sans racine carrée au dénominateur, puis simplifier.

\begin{tasks}[label=\textbf{\alph*.}, item-indent=2em, after-item-skip=0.6em](4)
	\task $\dfrac{1}{\sqrt{5}}$
	\task $\dfrac{10}{\sqrt{2}}$
	\task $\dfrac{3}{2\sqrt{3}}$
	\task $\dfrac{15}{\sqrt{45}}$
\end{tasks}
\end{EXO}

\begin{prop}{4. Racine d'une somme}
\rub{Propriété}Si $a > 0$ et $b > 0$ : \ $\sqrt{a + b} < \sqrt{a} + \sqrt{b}$. \ On ne « coupe » \textbf{jamais} la racine d'une somme.

\rub{Exemple}Comparer $\sqrt{9 + 16}$ et $\sqrt{9} + \sqrt{16}$.
\afaire[1.6em]{$\sqrt{9 + 16} = \sqrt{25} = 5$, \ mais \ $\sqrt{9} + \sqrt{16} = 3 + 4 = 7$.}

\rub{Démonstration}De deux nombres positifs, le plus grand a le plus grand carré.
\afaire[5.4em]{$\left(\sqrt{a} + \sqrt{b}\right)^2 = \left(\sqrt{a}\right)^2 + 2\sqrt{a}\sqrt{b} + \left(\sqrt{b}\right)^2 = a + b + 2\sqrt{a}\sqrt{b}$ \quad (identité remarquable).

$\left(\sqrt{a + b}\right)^2 = a + b$. \quad Comme $2\sqrt{a}\sqrt{b} > 0$, on a $\left(\sqrt{a} + \sqrt{b}\right)^2 > \left(\sqrt{a + b}\right)^2$.

Les deux nombres sont positifs, donc $\sqrt{a} + \sqrt{b} > \sqrt{a + b}$.}
\end{prop}

\begin{EXO}{}{}
\begin{enumerate}
	\item Léo écrit : $\sqrt{36 + 64} = 6 + 8 = 14$. Trouver son erreur et donner le bon résultat.
	\item Comparer, sans calculatrice : \quad \textbf{a.} $\sqrt{2} + \sqrt{8}$ et $\sqrt{10}$ \qquad \textbf{b.} $\sqrt{18} + \sqrt{32}$ et $\sqrt{98}$.
\end{enumerate}
\end{EXO}

% Corrigé : version prof uniquement (compile.sh produit aussi la version _ELEVE)
\ifdefined\ELEVE\else
\clearpage

\begin{Correction}
\small
\begin{EXO}{}{}
\begin{tasks}[label=\textbf{\alph*.}, item-indent=2em](2)
	\task $\res{0{,}3}$ \enspace ($0{,}3^2 = 0{,}09$)
	\task $\res{1\,000}$ \enspace ($\left(10^3\right)^2 = 10^6$)
	\task $\res{3}$
	\task $\sqrt{64} = \res{8}$
\end{tasks}
\end{EXO}

\begin{EXO}{}{}
\begin{tasks}[label=\textbf{\alph*.}, item-indent=2em](2)
	\task $\sqrt{36} = \res{6}$
	\task $\sqrt{225} = \res{15}$
	\task $9 \times 2 = \res{18}$
	\task $10\sqrt{18} = 10 \times 3\sqrt{2} = \res{30\sqrt{2}}$
\end{tasks}
\end{EXO}

\begin{EXO}{}{}
\begin{tasks}[label=\textbf{\alph*.}, item-indent=2em](3)
	\task $\sqrt{16\times3} = \res{4\sqrt{3}}$
	\task $\sqrt{25\times3} = \res{5\sqrt{3}}$
	\task $\sqrt{49\times2} = \res{7\sqrt{2}}$
	\task $\sqrt{100\times5} = \res{10\sqrt{5}}$
	\task $\sqrt{144\times2} = \res{12\sqrt{2}}$
\end{tasks}
\end{EXO}

\begin{EXO}{}{}
\begin{tasks}[label=\textbf{\alph*.}, item-indent=2em](2)
	\task $A = 10\sqrt{2} - 6\sqrt{2} + 3\sqrt{2} = \res{7\sqrt{2}}$
	\task $B = 5\sqrt{3} - 4\sqrt{3} + \sqrt{3} = \res{2\sqrt{3}}$
	\task $C = 6\sqrt{5} - 3\sqrt{5} + 4\sqrt{5} = \res{7\sqrt{5}}$
	\task $D = \sqrt{12} + \sqrt{27} = 2\sqrt{3} + 3\sqrt{3} = \res{5\sqrt{3}}$
\end{tasks}
\end{EXO}


\begin{EXO}{}{}
\begin{tasks}[label=\textbf{\alph*.}, item-indent=2em](4)
	\task $\res{\dfrac{5}{7}}$
	\task $\sqrt{9} = \res{3}$
	\task $\sqrt{\dfrac{4}{9}} = \res{\dfrac{2}{3}}$
	\task $\sqrt{\dfrac{4}{25}} = \res{\dfrac{2}{5}}$
\end{tasks}
\begin{tasks}[label=\textbf{\alph*.}, item-indent=2em, resume](3)
	\task $\dfrac{5\sqrt{2} + 2\sqrt{2}}{\sqrt{2}} = \dfrac{7\sqrt{2}}{\sqrt{2}} = \res{7}$
	\task $\dfrac{2\sqrt{3} - \sqrt{3}}{3\sqrt{3}} = \dfrac{\sqrt{3}}{3\sqrt{3}} = \res{\dfrac{1}{3}}$
	\task $\dfrac{2 \times 2\sqrt{5}}{3\sqrt{5}} = \res{\dfrac{4}{3}}$
\end{tasks}
\end{EXO}

\begin{EXO}{}{}
\begin{tasks}[label=\textbf{\alph*.}, item-indent=2em](4)
	\task $\res{\dfrac{\sqrt{5}}{5}}$
	\task $\dfrac{10\sqrt{2}}{2} = \res{5\sqrt{2}}$
	\task $\dfrac{3\sqrt{3}}{6} = \res{\dfrac{\sqrt{3}}{2}}$
	\task $\dfrac{15}{3\sqrt{5}} = \dfrac{5}{\sqrt{5}} = \res{\sqrt{5}}$
\end{tasks}
\end{EXO}



\begin{EXO}{}{}
\begin{enumerate}
	\item Il a « coupé » la racine d'une somme : $\sqrt{36 + 64} = \sqrt{100} = \res{10}$.
	\item \textbf{a.} D'après la propriété, \res{\sqrt{2} + \sqrt{8} > \sqrt{10}} \enspace (vérification : $\sqrt{2} + \sqrt{8} = 3\sqrt{2} = \sqrt{18}$).

	\textbf{b.} $\sqrt{18} + \sqrt{32} = 3\sqrt{2} + 4\sqrt{2} = 7\sqrt{2} = \sqrt{98}$ : \res{\text{les deux nombres sont égaux}}.
\end{enumerate}
\end{EXO}

\clearpage
\end{Correction}
\fi
\end{document}
